\section{The actual third growth target on the moving second hold}
\label{tge:section}

We continue from the actual second returned state, with its first
vorticity strictly positive. The third wave sees both older temperature
gradients and the changing first shear. This section proves attainment
of its original growth threshold, with the original third insertion
rule and retained positive physical diffusion. The second parent is
held by its original feedback during this interval. The third
transition changes that control and requires its own calculation;
no third return is inferred from the growth result.

\subsection{A further explicit choice in the same source family}

Retain every original datum, profile, frequency and constant from
Section~\ref{ave:section}. In particular $\nu=\sqrt{A_0}$,
$\mu=\lambda_0=\lceil\sqrt{A_0}/2\rceil$ and
$K_j=\mu\widehat\lambda_j$ are still physical quantities.
Write $\sigma_1=\sigma$ for the actual second-stage scale in
\eqref{ave:actual_scale}. Extend the exact source recursion to
\begin{equation}\label{tge:source}
\begin{gathered}
 Q_3=Q_*+3,\qquad
 \widehat\lambda_3=\widehat\lambda_2^{Q_3},\qquad
 k_3=\max\{k\in\mathbb N_0:120k\leq Q_3,
                         \ K(k)\leq\log\widehat\lambda_3\},\\
 L_3=(k_3+41/8)\log\widehat\lambda_3,\qquad
 S_3=(A_0/\mu)\widehat\lambda_3^{-k_3-6},\qquad
 P_3^*=(A_0/\mu)\widehat\lambda_3^{-7/8}=S_3e^{L_3}.
\end{gathered}
\end{equation}
The finite maximum exists for the same reason as the earlier two.
Define fixed numbers
\begin{gather*}
 c_2=(41/8)Q_2,\qquad c_3=(41/8)Q_2Q_3,\\
 C_3=(\lfloor Q_3/120\rfloor+41/8)Q_2Q_3,\\
 C_2=e^9\sqrt{5\sqrt{10}/4},\qquad
 C_M=\frac{256C_2}{15\sqrt e}.
\end{gather*}
Then $L_2\geq c_2\log\widehat\lambda_1$,
$L_3\geq c_3\log\widehat\lambda_1$, and
$L_3\leq C_3\log\widehat\lambda_1$ for all these exact choices.
Using the already defined $Y_0,C_\sigma,C_N$, select
\begin{equation}\label{tge:Y3}
 Y_3=1+\max\left\{\begin{gathered}
 Y_0,\quad 1024/c_2,\quad 16384/c_3,\quad
 \frac8{Q_2-1}\log(12C_\sigma^2),\\
 \frac{16}{Q_2-2}\log(2C_\sigma^2C_3/c_2),\quad
 \frac{\log(290C_N)}{3Q_2-49/16},\quad
 \frac{\log(145\sqrt{10}/s_0)}{2Q_2},\\
 \frac{\log(4C_NC_M)}{47Q_2/16-3},\quad
 \frac{\log(\sqrt{11}C_M/s_0)}{(31Q_2+1)/16}
 \end{gathered}\right\}.
\end{equation}
Choose $\widehat\lambda_1\geq e^{Y_3}$.
All denominators are positive because $Q_2\geq202$.
These additional lower bounds retain the same source frequency family
and integer selection; they do not replace any earlier scale by one.
Set
\begin{equation}\label{tge:d3}
 \sigma_- =\sqrt e\,\nu\widehat\lambda_1^{1/16},\qquad
 d_{3,\max}=\min\left\{d_{\max},
      \frac{\sigma_-\log2}{160K_2^2},
      \frac{L_3\sigma_-}{128K_3^2}\right\}>0.
\end{equation}
For the remainder choose any $0<d\leq d_{3,\max}$ and construct
the actual first and second stages already proved. Every number on
the right of \eqref{tge:d3} is specified before selecting the second
root or solving a new growth equation. The lower bound
$\sigma_1\geq\sigma_-$ follows from \eqref{ave:scale_bounds}.

\subsection{The actual third insertion scale}

Let $t_b=t_2^b$ be the selected second return, and use a subscript
$b$ for evaluation there. Set $a=\sin s_2$, $h_b=h(t_b)$,
$\chi_b=\sqrt{h_b^2+a^2}$, and $P_{2,b}=-\Theta_2(t_b)>0$.
Define the original second gradient amplitude and next scale at this
actual endpoint by
\begin{equation}\label{tge:birthscale}
 A_{2,b}=K_2\chi_bP_{2,b},\qquad
 \sigma_2^2=|G_{<3}(t_b)|,\qquad
 s_3=\frac{L_3\sigma_1}{\sigma_2},\qquad
 \Gamma_3=\sigma_2\sin s_3.
\end{equation}
In particular $\sigma_2$ is evaluated after the selected second
control, rather than at an inviscid endpoint.
The complete older gradient at this time is
\begin{gather*}
 G_{<3}(t_b)=-A_0R_{\alpha_b}e_2
              -A_{1,b}\zeta_{1,b}-A_{2,b}e_2,\\
 A_{1,b}=-K_1\Theta_1(t_b),\quad
 \zeta_{1,b}=\frac{(-a,h_b)}{\chi_b},\quad
 \alpha_b=s_*+\arctan(a/h_b).
\end{gather*}
Theorem~\ref{sr:theorem} gives
$\chi_b\geq1$, $P_{2,b}\geq P_2^*$,
$A_{1,b}\leq2\sigma_1^2$, and $A_0\leq\sigma_1^2$.
The new choice $L_2>1024$ and \eqref{ave:angle_bounds} imply
\begin{equation}\label{tge:a_small}
 0<a<\frac1{12HL_2}\leq\frac1{12288}.
\end{equation}
Because $h_b\geq\cos s_2\geq15/16$,
$0<\alpha_b\leq1/8+(16/15)a<1/4$.
Both older terms in the displayed $G_{<3}$ have negative vertical
component. Its norm is therefore at least $A_{2,b}$, while the
triangle inequality bounds it above by $A_{2,b}+3\sigma_1^2$.
Moreover \eqref{tge:Y3} and the original threshold give
\[
 A_{2,b}\geq A_0\widehat\lambda_2^{1/8},\qquad
 \frac{A_{2,b}}{\sigma_1^2}
 \geq\frac{\widehat\lambda_1^{(Q_2-1)/8}}{C_\sigma^2}>12.
\]
Consequently
\begin{equation}\label{tge:scale_bounds}
 A_{2,b}\leq\sigma_2^2\leq\frac54A_{2,b},\qquad
 \sigma_2>\sigma_1,\qquad
 \sigma_2\geq\nu\widehat\lambda_2^{1/16},\qquad
 \frac45\leq\frac{A_{2,b}}{\sigma_2^2}\leq1.
\end{equation}

The actual third angle obeys an inequality stronger than a branch
choice. Since $a\geq L_2\nu/(2\sigma_1)$, the scale bounds imply
\[
 \frac{s_3}{a}
 \leq\frac{2L_3\sigma_1^2}{L_2\nu\sigma_2}
 \leq2C_\sigma^2\frac{C_3}{c_2}
             \widehat\lambda_1^{-(Q_2-2)/16}<1.
\]
Thus $0<s_3<a<1/8$. The elementary inequality
$\sin x\geq x-x^3/6\geq(15/16)x$ for $0\leq x\leq1/8$
follows by integrating $\cos x\geq1-x^2/2$.
It proves
\begin{equation}\label{tge:rate_scale}
 \frac{15}{16}L_3\sigma_1\leq\Gamma_3\leq L_3\sigma_1.
\end{equation}
No small-angle identity has replaced the actual sine in the controls.

\subsection{Exact moving geometry and the changing coupling numerator}

Use physical elapsed time $r=t-t_b$ and the original third growth
control $\dot\alpha=F_2^{\rm ctl}$. The second amplitude and its
horizontal phase component remain zero in vorticity and phase,
respectively, while both older temperatures and the first shear evolve:
\begin{align}
 \zeta_1&=(-a,h)/\chi,&\zeta_2&=(0,\chi),&
 \chi&=\sqrt{h^2+a^2},&\dot h&=a\Omega_1,\nonumber\\
 \dot A_1&=-A_0\sin s_*\,\Omega_1-dK_1^2A_1,&
 \dot P_2&=-dK_2^2\chi^2P_2,&
 \dot\alpha&=-a^2\Omega_1/\chi^2.
 \label{tge:older_hold}
\end{align}
These are the previously proved holding equations in physical time.
Let
\begin{equation}\label{tge:detconstants}
 \beta_b=\arctan(a/h_b),\quad
 b=\sin(s_3+\beta_b),\quad g_b=\cos(s_3+\beta_b),\quad
 c=bh_b-ag_b=\chi_b\sin s_3>0.
\end{equation}
Here $b$ is an oriented determinant magnitude, not an amplitude
coefficient; the latter will carry its layer index below.
Since $s_3\leq a$ and $\beta_b\leq(16/15)a$,
$0<s_3+\beta_b<3a<1/8$, so $g_b>0$ and $0<b<3a$.

The third material matrix and phase have the exact expressions
\begin{align}
 M_3(t)&=M_2(t)M_2(t_b)^{-1},& M_3(t_b)&=I,\nonumber\\
 g(t)&=g_b+(b/a)(h(t)-h_b),&
 \zeta_3(t)&=g\zeta_1-bJ\zeta_1
       =\frac{(c,hg+ab)}{\chi},& R_3&=|\zeta_3|^2=g^2+b^2.
 \label{tge:phase}
\end{align}
Here is a direct verification preserving the original datum
$p_3=e(s_3)$. Since $\Omega_2=0$, the actual matrices
$D_{<3}$ and $D_{<2}$ are equal, so differentiation proves the
matrix equation for $M_3$ with its stated initial value.
In the basis $\zeta_1,J\zeta_1$, the third phase equation is
$\dot g=b\Omega_1$, $\dot b=0$; these follow by differentiating
its scalar products and using
$\dot\zeta_1=\dot\alpha J\zeta_1$,
$\dot\zeta_3=\dot\alpha J\zeta_3
 -\Omega_1\zeta_1(J\zeta_1\cdot\zeta_3)$.
Integration and $\dot h=a\Omega_1$ give \eqref{tge:phase}.
At birth, resolving $e(s_3)$ in this oriented basis gives precisely
$g_b,b$ in \eqref{tge:detconstants}. The identity for $c$ follows
from $h_b=\chi_b\cos\beta_b$, $a=\chi_b\sin\beta_b$ and
the sine subtraction formula. This verifies both components of the
phase formula and its initial datum. Uniqueness of the linear phase
equation identifies it with $M_3^{-T}p_3$.

The full inherited gradient and third amplitude coefficients are
\begin{align}
 G_{<3}&=-(A_1+A_0\cos s_*+K_2P_2h)\zeta_1
              +(-A_0\sin s_*+aK_2P_2)J\zeta_1,\nonumber\\
 N_3&=b(A_1+A_0\cos s_*)+gA_0\sin s_*+cK_2P_2,\nonumber\\
 a_3&=\frac{N_3}{K_3R_3},\qquad
 b_3=K_3\zeta_{3,1}=\frac{K_3c}{\chi}.
 \label{tge:coefficients}
\end{align}
To verify the first line, use
$R_\alpha e_2=\cos s_*\zeta_1+\sin s_*J\zeta_1$ and
$\zeta_2=h\zeta_1-aJ\zeta_1$ in the original sum
$G_{<3}=-A_0R_\alpha e_2-A_1\zeta_1-K_2P_2\zeta_2$.
The scalar product of $J\zeta_3=gJ\zeta_1+b\zeta_1$ with
that sum is $-N_3$, proving the coefficient including the original
gravity direction. Differentiating its full numerator gives
\begin{equation}\label{tge:Ndot}
 \dot N_3=-d\{bK_1^2A_1+cK_2^3\chi^2P_2\}.
\end{equation}
Indeed the terms $-bA_0\sin s_*\Omega_1$ from $\dot A_1$
and $b\Omega_1 A_0\sin s_*$ from $\dot g$ cancel exactly;
the second temperature contributes the other displayed term.
Thus the inviscid numerator would be exactly constant despite the
changing shear. At positive $d$ both parent losses are retained.

\subsection{A common proved horizon before taking an amplitude quotient}

Set $R_H=16/\sigma_1$. Its duration in the second-stage time is
$\Gamma_2R_H=16a<1/32$, by \eqref{tge:a_small}.
Thus the entire closed older system exists on $0\leq r\leq R_H$
by its already proved second hold, uniformly for the current $d$.
There $A_1,P_2$ are positive, $\Omega_1<9\sigma_1$,
$\chi\geq\chi_b\geq1$, and
\begin{equation}\label{tge:horizon_geometry}
\begin{gathered}
 0\leq h-h_b\leq144a,\qquad
 0\leq g-g_b\leq144b,\\
 1\leq\sqrt{R_3}\leq1+432a\leq265/256<\sqrt2,
 \qquad1\leq\frac{\chi}{\chi_b}\leq259/256<2.
\end{gathered}
\end{equation}
The first inequality integrates $\dot h=a\Omega_1$.
The second uses the exact relation \eqref{tge:phase}.
For the third apply the Euclidean triangle inequality to
$(g,b)=(g_b,b)+(g-g_b,0)$, whose first summand has norm one.
The fourth follows similarly from
$\sqrt{h^2+a^2}\leq\chi_b+(h-h_b)$ and $\chi_b\geq1$.
The rational constants use $432/12288=9/256$ and
$144/12288=3/256$; $265^2<2\cdot256^2$ proves the strict square-root
bound without a numerical approximation.

All four terms in $N_3$ are positive on this horizon.
Since $K_1<K_2$ and $\chi^2<10$, \eqref{tge:Ndot} gives
$-10dK_2^2N_3\leq\dot N_3\leq0$.
Multiplication by $e^{10dK_2^2r}$ and integration, followed by
\eqref{tge:d3}, proves
\begin{equation}\label{tge:coeffbox}
 \frac12N_i\leq N_3(r)\leq N_i,\qquad
 \frac14a_i\leq a_3(r)\leq a_i,\qquad
 \frac12b_i\leq b_3(r)\leq b_i,
\end{equation}
where $N_i=N_3(t_b)$, $a_i=N_i/K_3$, and
$b_i=K_3\sin s_3$ are the actual birth values.
Indeed $10dK_2^2R_H=160dK_2^2/\sigma_1\leq\log2$;
the remaining bounds use \eqref{tge:horizon_geometry} in the exact
coefficient formulas. No value of a later temperature is assumed
in these estimates.

For the actual initial rate $\gamma_i=\sqrt{a_ib_i}$ put
$n_i=\gamma_i^2/\Gamma_3^2=N_i/(\sigma_2^2\sin s_3)$.
The positive newest-parent contribution to its numerator is
$cK_2P_{2,b}=A_{2,b}\sin s_3$.
Its remaining two terms satisfy
\[
 b(A_{1,b}+A_0\cos s_*)+g_bA_0\sin s_*\leq4\sigma_1^2,
\]
using $b,g_b\leq1$, $A_{1,b}\leq2\sigma_1^2$ and
$A_0\leq\sigma_1^2$.
Equations~\eqref{tge:scale_bounds} and \eqref{tge:rate_scale}
therefore give
\begin{equation}\label{tge:initial_rate}
 \frac45\leq n_i\leq1+\frac{64\sigma_1}{15L_3\sigma_2}
       \leq\frac{16}{15},\qquad
 \gamma_i\geq\frac78\Gamma_3
       \geq\frac{105}{128}L_3\sigma_1.
\end{equation}
For the middle upper bound use $L_3\geq64$ and $\sigma_2>\sigma_1$.
The lower rate follows from $4/5>(7/8)^2$.

The original third seed is exactly
$\Theta_3(t_b)=-S_3$, $\Omega_3(t_b)=-S_3\sqrt{b_i/a_i}$.
Solve its complete linear system with the changing coefficients:
\begin{equation}\label{tge:pair}
 \dot\Theta_3=a_3\Omega_3-dK_3^2R_3\Theta_3,
 \qquad\dot\Omega_3=b_3\Theta_3-dK_3^2R_3\Omega_3.
\end{equation}
Its coefficients and derivatives are finite on the common horizon,
so the iterated integral series provides existence and uniqueness
there. Define
\[
 E(r)=\exp\left(dK_3^2\int_0^r R_3(t_b+u)\,du\right),\qquad
 X=-E\Theta_3,\qquad Y=-E\Omega_3.
\]
Differentiation gives $\dot X=a_3Y$, $\dot Y=b_3X$.
Both initial values are positive, and before a hypothetical first zero
both derivatives are positive. This excludes the zero and proves
$\Theta_3,\Omega_3<0$ throughout the horizon. The quotient may now
be formed. With $u=\Omega_3/\Theta_3$ and $u_i=\sqrt{b_i/a_i}$,
it satisfies $\dot u=b_3-a_3u^2$. The bounds in
\eqref{tge:coeffbox} make $u_i/\sqrt2$ a constant subsolution and
$2u_i$ a constant supersolution. The integrating-factor comparison
for their differences proves
\begin{equation}\label{tge:quotient}
 \frac{u_i}{\sqrt2}\leq u\leq2u_i,
 \qquad \sqrt{15/32}\leq
       v_3:=\frac{\sigma_2\Omega_3}{K_3\Theta_3}\leq\sqrt5<3.
\end{equation}
The exact identity $\sigma_2u_i/K_3=1/\sqrt{n_i}$ and
\eqref{tge:initial_rate} give the latter interval.

\subsection{The original third threshold and its full diffusion cost}

The last bound in \eqref{tge:d3} and
\eqref{tge:initial_rate} imply
$dK_3^2\leq L_3\sigma_1/64\leq\gamma_i/(16\sqrt2)$.
For the last comparison it suffices to use
$\gamma_i\geq L_3\sigma_1/2$ and $\sqrt2\leq2$.
Thus the actual $P_3=-\Theta_3$ has the positive physical rate
\begin{equation}\label{tge:gain_rate}
 \frac{\gamma_i}{8\sqrt2}
 \leq(\log P_3)'=a_3u-dK_3^2R_3
 \leq2\gamma_i<3\Gamma_3.
\end{equation}
The lower bound uses $a_3u\geq a_iu_i/(4\sqrt2)$,
$R_3<2$, and the displayed diffusion bound. The upper bound
uses $a_3u\leq2a_iu_i$ and $n_i\leq16/15<9/4$.

\begin{theorem}\label{tge:actual_threshold}
For every fixed original datum and profiles, choose the source
frequencies by \eqref{tge:Y3} and any $0<d\leq d_{3,\max}$ in
\eqref{tge:d3}. After the actual two returns already constructed,
the original third growth control reaches its original target
$P_3^*$ exactly once, at $t_3^a=t_b+r_a$, with
\begin{equation}\label{tge:times}
 \frac{L_3}{2\gamma_i}\leq r_a
       \leq\frac{8\sqrt2L_3}{\gamma_i}
       \leq\frac{1024\sqrt2}{105\sigma_1}
       <\frac{14}{\sigma_1}<\frac{16}{\sigma_1}.
\end{equation}
The crossing is transverse, its activation is already complete,
and its initial and final physical amplitude ratios obey
\eqref{tge:quotient}. Both inherited temperatures and the first
shear evolve by \eqref{tge:older_hold} throughout.
The exact logarithmic diffusion cost of the third amplitude is
\begin{equation}\label{tge:diffusion_cost}
 dK_3^2\int_0^{r_a}R_3(t_b+r)\,dr,
 \qquad
 dK_3^2r_a\leq
 dK_3^2\int_0^{r_a}R_3(t_b+r)\,dr
                  \leq2dK_3^2r_a.
\end{equation}
Its total actual logarithmic gain is exactly $L_3$.
\end{theorem}
\begin{proof}
The complete older and new systems were constructed on a fixed horizon
before the threshold was chosen. The upper candidate time in
\eqref{tge:times} is within that horizon: squaring positive numbers
gives $(1024\sqrt2)^2<1470^2$. Integrating the lower rate in
\eqref{tge:gain_rate} to that candidate time gives at least $L_3$.
Continuity and strict increase therefore prove a unique crossing.
Integrating the upper rate gives its lower time bound.
Because $2\gamma_i<3\Gamma_3$ and $L_3\geq64$,
$\Gamma_3r_a>L_3/3>1$, so the original activation
$\eta(\Gamma_3r)$ is one before this crossing.
The rate lower bound proves transversality. The positive phase-length
bounds in \eqref{tge:horizon_geometry} give
\eqref{tge:diffusion_cost} by integration, with all physical factors
retained. The target identity in \eqref{tge:source} gives the exact
gain $L_3$.
\end{proof}

For the two older amplitudes during this third growth, the exact
diffusion exponents are $dK_1^2r_a$ and
$dK_2^2\int_0^{r_a}\chi(t_b+r)^2dr$. The second one is at most
$10dK_2^2r_a\leq\log2$ on our common horizon. The additional
first-temperature coupling $-A_0\sin s_*\Omega_1$ is retained in
\eqref{tge:older_hold}; it is not a damping exponent.
The numerator's exact total change is the integral of
\eqref{tge:Ndot}, including both parent contributions.

\paragraph{Exact third support nesting.}\label{tge:nesting}
The source third support also fits the changing exact affine cells.
Indeed the relative matrix \eqref{tge:phase} is a product of two
rotations and the shear with off-diagonal entry $-(h-h_b)/a$;
hence $\|M_3\|,\|M_3^{-1}\|\leq145$ on this horizon.
The older bound is $\|M_2^{-1}\|\leq N\leq
C_N\widehat\lambda_1^{1/16}$. With the exact source radius
$\ell_3=\widehat\lambda_2^{-3}/\mu$, the two growth-support conditions
in \eqref{tge:Y3} give
\[
 2\cdot145N\widehat\lambda_2^{-3}<\widehat\lambda_1^{-3},
 \qquad145\sqrt{10}\widehat\lambda_2^{-2}<s_0.
\]
Thus $|M_2^{-1}x|<\ell_2/2$ and
$|K_2\zeta_2\cdot x|<s_0$ on the entire third support.
The earlier proved nesting places it in the first and base affine
cells as well. These explicit inequalities justify the complete
compact third spatial increment and its full forces below.

This proves a further original growth event with retained diffusion,
not a third return or an infinite iteration. The newly permitted
range includes a factor $K_3^{-2}$ and retains that dependence.
The following transition must evolve both older vorticities and its
prescribed blend of $F_2^{\rm ctl}$ and $F_3^{\rm ctl}$; the exact
third-growth numerator identity identifies the parent losses that
must remain in that calculation.
